% !PRINTED 597
\weberpage{597}

Die Gesamtheit der ganzen und gebrochenen Zahlen des Körpers $\Omega$, der wir
einen beliebigen Exkludenten $S$ beilegen, bildet die Zahlgruppe $O$ und die
Gesamtheit der Ideale $\overline{O}$ in dem gleichen Sinne eine \gesperrt{Idealgruppe}.

Ebenso bildet die Gesamtheit der \gesperrt{numerischen Einheiten} eine Zahlgruppe $E$
und die Gesamtheit der \gesperrt{funktionalen Einheiten} eine Idealgruppe $\overline{E}$.

Die Multiplikation zweier Gruppen $A$ und $B$ zu dem Produkt $AB$ geschieht
dadurch, daß man jedes Element von $A$ mit jedem Element von $B$
multipliziert. Ist $A$ eine Zahlgruppe, so ist $\overline{E}A$ eine
Idealgruppe. Diese enthält aber nur Hauptideale und kann daher
\gesperrt{Hauptidealgruppe} genannt werden.

Wenn die ganze Gruppe $O$ in eine \gesperrt{endliche Anzahl von Nebengruppen} nach $A$
zerfällt:
\weq{1}{O = A\alpha_1 + A\alpha_2 + \cdots + A\alpha_j,}
so heißt die Zahl $j$ der Index von $A$ und wird mit
\weq{2}{j = (O, A)}
bezeichnet (wie in § 100).

Wir sprechen den ersten Satz aus:

1. \gesperrt{Ist $A$ eine Zahlgruppe mit endlichem Index $j$, und $\eta$ eine
beliebige Zahl in $O$, so ist $\eta^j$ in $A$ enthalten.}

Denn von der unbegrenzten Reihe der Potenzen $1, \eta, \eta^2, \eta^3, \dots$,
müssen zwei verschiedene Glieder in derselben Nebengruppe $A\alpha_i$
enthalten sein, und ihr Quotient, der ja auch eine Potenz von $\eta$ ist, ist
in $A$ enthalten. Der niedrigste Exponent $r$, für den $\eta^r$ in $A$
enthalten ist, erweist sich in bekannter Weise als Teiler von $j$.

2. \gesperrt{Ist $A$ eine Zahlgruppe mit endlichem Index, so läßt sich auch
die Idealgruppe $\overline{O}$ nach $\overline{E}A$ in eine endliche Zahl von
Nebengruppen zerlegen:}
\weq{3}{\begin{aligned}
\overline{O} &= \overline{A}_1 + \overline{A}_2 + \cdots + \overline{A}_h,\\
&= A\mathfrak{a}_1 + A\mathfrak{a}_2 + \cdots + A\mathfrak{a}_h,
\end{aligned}}
und zwar können wir darin $\mathfrak{a}_1, \mathfrak{a}_2, \dots,
\mathfrak{a}_h$ als ganze Ideale in $\overline{O}$ wählen, da wir jedes
gebrochene Ideal in $\overline{O}$ nach 1. durch Multiplikation mit einer
Zahl in $A$ in ein ganzes Ideal verwandeln können. Wir brauchen nur für
$\eta$ eine Zahl zu wählen, die durch den Nenner des Ideals teilbar ist, und
mit einer genügend hohen Potenz von $\eta$ zu multiplizieren.
